% --- SECTION 2: THEORETICAL BACKGROUND ---
\section{Theoretical Background}

\subsection{Reinforcement Learning Framework}

The trading problem is formulated as a Markov Decision Process (MDP), defined by the tuple 
$(\mathcal{S}, \mathcal{A}, P, R, \gamma)$, where $\mathcal{S}$ denotes the state space, 
$\mathcal{A}$ the action space, $P$ the transition dynamics, $R$ the reward function, and 
$\gamma \in (0,1)$ the discount factor.

At each time step $t$, the agent observes a state $s_t$ (constructed from market features), 
selects an action $a_t$ according to a parameterized stochastic policy $\pi_\theta(a|s)$, 
receives reward $r_{t+1}$, and transitions to the next state $s_{t+1}$.

The reward function is defined as the logarithmic return of the portfolio valuation:
\[
r_t = \ln\left(\frac{p_t}{p_{t-1}}\right),
\]
where $p_t$ denotes the portfolio valuation at time step $t$.

The objective is to maximize the expected discounted return:

\[
J(\theta) = \mathbb{E}_{\pi_\theta} 
\left[ 
\sum_{t=0}^{T} \gamma^t r_{t+1} 
\right].
\]

In contrast to value-based methods, policy gradient approaches directly optimize 
the policy parameters $\theta$ to increase expected return \citep{wikipediapolicysradient}.

\subsection{Monte Carlo Policy Gradient (REINFORCE)}

Monte Carlo methods estimate returns using complete sampled episodes. 
For a trajectory of length $T$, the return from time step $t$ is defined as:

\[
G_t = \sum_{k=0}^{T-t} \gamma^k r_{t+k+1}.
\]

Unlike temporal-difference methods, Monte Carlo approaches do not use bootstrapping; 
they rely solely on actual observed rewards. This leads to unbiased return estimates, 
but potentially high variance.

Using the log-derivative trick, the gradient of the objective function can be written as:

\[
\nabla_\theta J(\theta) =
\mathbb{E}_{\pi_\theta}
\left[
G_t \nabla_\theta \log \pi_\theta(a_t|s_t)
\right].
\]

This result leads to the REINFORCE update rule \citep{williams1992simple}:

\[
\theta \leftarrow \theta 
+ \alpha \, G_t \nabla_\theta \log \pi_\theta(a_t|s_t),
\]

where $\alpha$ is the learning rate.

In practice, the expectation is approximated using sampled trajectories. 
Updates are performed after each episode, making the algorithm fully Monte Carlo.

\subsection{Variance and Baseline Reduction}

A key limitation of REINFORCE is the high variance of its gradient estimator. 
In financial environments, rewards are noisy and highly volatile, which can 
lead to unstable learning and slow convergence.

To reduce variance without introducing bias, a baseline function $b(s_t)$ 
can be subtracted from the return:

\[
\nabla_\theta J(\theta)
=
\mathbb{E}
\left[
(G_t - b(s_t)) \nabla_\theta \log \pi_\theta(a_t|s_t)
\right].
\]

If the baseline depends only on the state and not on the action, 
the estimator remains unbiased. A common choice is the state-value function 
$V(s_t)$, which estimates the expected return from state $s_t$.

Defining the advantage as

\[
A_t = G_t - V(s_t),
\]

the update becomes:

\[
\theta \leftarrow \theta 
+ \alpha \, A_t \nabla_\theta \log \pi_\theta(a_t|s_t).
\]

The subtraction centers the gradient signal around zero, 
reducing variance and improving stability. 

In this project, the performance of standard REINFORCE 
is compared to REINFORCE with a learned value-function baseline. 
The comparison focuses on convergence behavior and training stability 
in a noisy financial environment.


\subsection{Practical Stabilization Techniques}
\label{sec:stabilization}

To address instability and high variance in neural network-based reinforcement learning \citep{wikipediapolicysradient}, these techniques were implemented:

\begin{enumerate}
    \item \textbf{Return Normalization:}
    \[
    \hat{G}_t = \frac{G_t - \mu_G}{\sigma_G + \epsilon}
    \]
    Normalizing returns reduces fluctuations in gradient scale, which helps stabilize neural network training, especially when episodic returns can be highly volatile.

    \item \textbf{Gradient Clipping:} \\
    To prevent exploding gradients during backpropagation, the norm of the gradient $g$ is clipped to a maximum value $c$:
    \[
    \|g\| \leq c
    \]
    This ensures stable policy updates.

    \item \textbf{Softmax Policy Parameterization:} \\
    The policy is typically parameterized using a softmax over action preferences to guarantee a valid probability distribution:
    \[
    \sum_a \pi_\theta(a|s) = 1
    \]
    This preserves sufficient exploration and prevents the policy from collapsing to deterministic actions too early.
\end{enumerate}

These stabilization techniques are essential for robust RL training in noisy financial domains.
